% !TEX program = xelatex
\documentclass{英语试卷}
\begin{document}
\examtitle{2025年合肥市高三第一次教学质量检测}{英语}
\begin{notice}
\par 1. 答卷前，考生务必将自己的姓名、准考证号填写在答题卡上。\par
\par 2. 回答选择题时，选出每小题答案后，用铅笔把答题卡上对应题目的答案标号涂黑。如需改动，用橡皮擦干净后，再选涂其他答案标号。回答非选择题时，将答案写在答题卡上，写在本试卷上无效。\par
\par 3. 考试结束后，将本试卷和答题卡一并交回，\par
\end{notice}
\parttitle{第一部分 听力（共两节，满分 30 分）}
\par 做题时，先将答案标在试卷上。录音内容结束后，你将有两分钟的时间将试卷上的答案转涂到答题卡上。\par
\sectitle{第一节（共5小题；每小题1.5分，满分7.5分）}
\par 听下面 5 段对话。每段对话后有一个小题，从题中所给的 A、B、C 三个选项中选出最佳选项，并标在试卷的相应位置。听完每段对话后，你都有 10 秒钟的时间来回答有关小题和阅读下一小题。每段对话仅读一遍。\par
\question{例：How much is the shirt?}
\begin{choices}(3)
  \item £ 19.15.
  \item £ 9.18.
  \item £ 9.15. 答案是 C。
\end{choices}
\question{1. What is the man doing?}
\begin{choices}(3)
  \item Applying for a job.
  \item Inquiring about a course.
  \item Registering for an activity.
\end{choices}
\question{2. What are the speakers talking about?}
\begin{choices}(3)
  \item A movie.
  \item A pianist.
  \item A website.
\end{choices}
\question{3. Where does the conversation probably take place?}
\begin{choices}(3)
  \item At a parking lot.
  \item At a train station.
  \item At a coffee house.
\end{choices}
\question{4. Why does the man make the phone call?}
\begin{choices}(3)
  \item To discuss a project.
  \item To ask about an exam.
  \item To get his stuff back.
\end{choices}
\question{5. How much does the man have to pay?}
\begin{choices}(3)
  \item 30 cents.
  \item 60 cents.
  \item 90 cents.
\end{choices}
\sectitle{第二节（共 15 小题；每小题 1.5 分，满分 22.5 分）}
\par 听下面 5 段对话或独白。每段对话或独白后有几个小题，从题中所给的 A、B、C 三个选项中选出最佳选项，并标在试卷的相应位置。听每段对话或独白前，你将有时间阅读各个小题，每小题 5 秒钟；听完后，各小题将给出 5 秒钟的作答时间。每段对话或独白读两遍。\par
\par 听第 6 段材料，回答第 6、7 题，\par
\question{6. Why is the man leaving his office?}
\begin{choices}(3)
  \item To pick up his son.
  \item To meet with a client.
  \item To handle an emergency.
\end{choices}
\question{7. What will the woman probably do next?}
\begin{choices}(3)
  \item Rearrange the meeting.
  \item Print the documents.
  \item Contact the driver.
\end{choices}
\par 听第 7 段材料，回答第 8 至 10 题。\par
\question{8. What does Bob do?}
\begin{choices}(3)
  \item A waiter.
  \item A chef.
  \item A grocer.
\end{choices}
\question{9. What will the woman's daughter have for the meal?}
\begin{choices}(3)
  \item Chicken Salad.
  \item Beef Steak.
  \item French Fries.
\end{choices}
\question{10. Why does the woman skip starters and desserts?}
\begin{choices}(3)
  \item She tries to lose some weight.
  \item She wants to avoid wasting food.
  \item She thinks they are overcharged.
\end{choices}
\par 听第 8 段材料，回答第 11 至 13 题。\par
\question{11. What is wrong with the jacket?}
\begin{choices}(3)
  \item It has a tear in it.
  \item It's out of fashion.
  \item It's the wrong size.
\end{choices}
\question{12. Why does the woman refuse to give back the money for the jacket?}
\begin{choices}(3)
  \item It has no proof of purchase.
  \item It's past the return period.
  \item It's been worn and washed.
\end{choices}
\question{13. How does the man feel about the service in the end?}
\begin{choices}(3)
  \item Satisfied.
  \item Annoyed.
  \item Confused.
\end{choices}
\par 听第 9 段材料，回答第 14 至 17 题。\par
\question{14. What does the woman do for the Tree-Planting Day?}
\begin{choices}(3)
  \item Set the place.
  \item Design a poster.
  \item Collect some information.
\end{choices}
\question{15. What can students benefit from the event according to the woman?}
\begin{choices}(3)
  \item Better stress management.
  \item Closer bonds with partners.
  \item More knowledge of nature.
\end{choices}
\question{16. What is the woman's attitude towards the follow-up activity?}
\begin{choices}(3)
  \item Supportive.
  \item Doubtful.
  \item Uninterested.
\end{choices}
\question{17. What will the man do this Saturday?}
\begin{choices}(3)
  \item Meet Eva.
  \item Plant trees.
  \item Buy tools.
\end{choices}
\par 听第 10 段材料，回答第 18 至 20 题，\par
\question{18. Who is the speaker probably talking to?}
\begin{choices}(3)
  \item Start-up founders.
  \item College students.
  \item Grocery suppliers.
\end{choices}
\question{19. What does "Hungry Harvest" do?}
\begin{choices}(3)
  \item Donate money to families in need.
  \item Supply fresh food to green grocers.
  \item Recover and sell imperfect produce.
\end{choices}
\question{20. What does the speaker expect the audience to do?}
\begin{choices}(3)
  \item Put in consistent effort.
  \item Focus on achievable goals.
  \item Take on social responsibility.
\end{choices}
\parttitle{第二部分 阅读理解（共两节，满分50分）}
\sectitle{第一节（共15小题；每小题2.5分，满分37.5分）}
\par 阅读下列短文，从每题所给的 A、B、C 和 D 四个选项中，选出最佳答案，\par
\par Visit These Cultural Hot Spots in 2024\par
\par Guangzhou Museum of Art, Guangzhou, China\par
\par Covered with 21,008 pieces of "fish scales(鳞)” glass, Guangzhou Museum of Art looks like a blooming hero flower. This landmark building features seven floors of exhibits including ancient and contemporary Chinese paintings, watercolors, calligraphy and sculpture.\par
\par Opening Hours: Tuesday to Sunday: 9 a. m-5 p. m. Closed: Monday\par
\par Istanbul Modern, Istanbul, Turkey\par
\par In Istanbul, there's a 110,000-square-foot art space given to Renzo Piano, a famous Italian architect, to show modern art from Turkey and beyond. The building has a shiny aluminum (铝) surface that mirrors the nearby Bosporus. Inside, glass-walled rooms feature large abstract paintings. The café at the top offers a breathtaking view of Asia, with a small "pool" that seems to stretch to the horizon.\par
\par Opening Hours: Tuesday to Thursday, Saturday and Sunday: 10 a. m.-6 p. m.\par
\par Friday: 10 a. m.-8 p. m.\par
\par Closed: Monday\par
\par Folger Shakespeare Library, Washington, D. C., USA\par
\par Fans of the Bard will find the world's largest collection of Shakespeareana, including writings or items related to Shakespeare, at this newly restored and expanded museum and library reopening in June. Located in Washington, D.C.'s Capitol Hill neighborhood, the Folger has more First Folios — the rare 1623 book containing all of Shakespeare's plays — than anywhere else in the world. An Elizabethan-style stage presents both classic and new works.\par
\par Opening hours: Tuesday, Wednesday and Sunday: 11 a. m.-6p. m. Thursday, Friday and Saturday: 11 a. m.-9 p. m.\par
\par Closed: Monday.\par
\par Canadian Canoe Museum, Ontario, Canada\par
\par One of the best things about visiting the world's largest collection of canoes(独木舟) and kayaks(皮艇)? You can actually row up to it. The streamlined, boat-shaped facility, which has its grand opening in May, is located lakeside in Peterborough. After seeing more than 600 boats or learning to carve paddles(桨) with Anishinaabe artists, visitors can go on a group boat trip around Little Lake.\par
\par Opening hours: Monday to Friday (May to October): 10 a. m.-5 p. m. Monday to Friday (November to April): 10 a. m.-4 p. m. Saturday and Sunday: 1 p. m.-4 p. m.\par
\question{21. What is a shared feature of Guangzhou Museum of Art and Istanbul Modern?}
\begin{choices}(2)
  \item They have unique architectural designs.
  \item They mainly display ancient art.
  \item They have exactly the same opening hours.
  \item They host some of the finest cafés.
\end{choices}
\question{22. What can visitors do at Folger Shakespeare Library?}
\begin{choices}(2)
  \item View a collection of modern art.
  \item Get copies of First Folios for free.
  \item Restore the collected writings and items.
  \item Enjoy plays on an Elizabethan-style stage.
\end{choices}
\question{23. What are the Sunday opening hours of the Canadian Canoe Museum in December?}
\begin{choices}(2)
  \item 10:00-16:00.
  \item 10:00-17:00.
  \item 11:00-21:00.
  \item 13:00-16:00.
\end{choices}
\passagelabel{B}
\par It was my birthday, and I was celebrating it as I had for years: Watching the annual Perseid meteor shower(流星雨). But this time, as the meteors flashed across the sky, I, a junior girl, began to question, "Could my lifelong love for astronomy be more than a hobby?"\par
\par I grew up in a small town, where even graduating from high school was not a given for many students. But I was driven to pursue my education, and like my best friend's dad — a pharmacist(药剂师), one of the few professional role models in my neighborhood, therefore, I chose to major in pharmacy, the science of preparing medicines, at the University of North Carolina.\par
\par The university was only a 3.5-hour drive away, but it felt across the universe. In the weeks since I began my pharmacy degree, I had met people with jobs and PhDs in different branches of science I had never considered — neuroscience, archaeology, and even astrophysics. Their passion and knowledge opened my eyes to worlds I had never imagined. As my birthday and the Perseid meteor shower approached, I was faced with a decision: Should I stick to my long-standing plans or take a step into the unknown?\par
\par Ultimately, a solar eclipse(日食) helped me make my decision. As the Moon gradually shaded the Sun, the world darkened around me. The stillness of the water, the desperate calls of birds, and the fading sunset deeply moved me. I knew what I had to do. I dropped out of pharmacy school to pursue previously unimagined options.\par
\par Facing years of challenges and growth, I worked odd jobs before earning a second degree in physics, starting a PhD in astronomy, and realized that the night sky's vastness inspired me to overcome doubts and follow my dreams. As I approach the final year of my PhD, I realize that growing up with the night sky gave me a valuable gift. It gave me an insight into life: On the scale of our vast universe, our individual needs, desires, fears, and insecurities shrink to insignificance, bringing humanity together. It liberated me from allowing earthly worries, shame, or self-doubt to stand in the way of pursuing new dreams.\par
\question{24. What motivated the author to major in pharmacy?}
\begin{choices}(2)
  \item Her ambition to become a doctor.
  \item The encouragement of her friend.
  \item Her childhood passion for pharmacy.
  \item The influence of her best friend's father.
\end{choices}
\question{25. How did the author's experience at university shape her?}
\begin{choices}(2)
  \item It helped her stick to her belief.
  \item It opened her mind to new ideas.
  \item It enhanced her pharmacy study.
  \item It developed her thinking abilities.
\end{choices}
\question{26. What can we learn from the author's understanding of life?}
\begin{choices}(2)
  \item Accepting insecurities helps us handle challenges.
  \item Personal problems still matter in the vast universe.
  \item It's vital to pursue dreams despite personal challenges.
  \item It is small and unimportant things that mean a lot to us.
\end{choices}
\question{27. Which of the following best describes the author?}
\begin{choices}(2)
  \item Caring and patient.
  \item Modest and honest.
  \item Generous and helpful.
  \item Curious and determined.
\end{choices}
\passagelabel{C}
\par Scientists with the BrainGate research center have developed a breakthrough brain-computer interface (BCI). The system aims to restore communication for individuals with speech disabilities caused by conditions like amyotrophic lateral sclerosis (ALS), which lead to the gradual loss of mobility and speech. It uses implanted sensors to interpret brain signals when a user attempts to speak, transforming them into text that is read aloud by a computer. The study, published in the New England Journal of Medicine on August 14, was led by David Brandman and Sergey Stavisky, faculty members at UC Davis Health.\par
\par ALS, also known as Lou Gehrig's disease, affects nerve cells controlling muscle movement. BCI technology can restore communication for those affected. In this study, Casey Harrell, a 45-year-old man with ALS, regained the ability to communicate effectively within minutes of using the system. Harrell described the experience as life-changing, saying, "My past condition was so demoralizing. Not being able to communicate seems as if you were trapped. Something like this technology will help people get back into life and society."\par
\par The BrainGate clinical trial, led by Dr. Leigh Hochberg from Brown University, has made major progress in interpreting brain signals to help people with paralysis(瘫痪) communicate. In July 2023, the UC Davis team implanted a BCI device into Harrell's brain, focusing on the left precentral gyrus, a part of the brain that controls speech. The system translates brain activity into phonemes (音素) and words, enabling real-time communication. Harrell's BCI generated speech with a voice similar to his own, using software trained with existing voice samples of his voice before he developed ALS. This allowed him to connect with his loved ones in a deeply personal way.\par
\par Over 32 weeks and 248 hours of use, the system achieved remarkable results, with 99.6\% accuracy for a 50-word vocabulary in 30 minutes of training and 90.2\% accuracy for a 125,000-word vocabulary after 1.4 additional training hours.\par
\par "This is transformative technology," Brandman said. "It provides hope for people who want to speak but can't." Dr. John Ngai, director of the NIH's BRAIN Initiative, highlighted the milestone as a major step in restoring communication for individuals locked in by paralysis.\par
\question{28. What do we know about the BrainGate BCI?}
\begin{choices}(2)
  \item It is used to treat speech disabilities.
  \item It translates brain signals into speech.
  \item It helps people with ALS move and speak.
  \item It connects brains to make computers smarter.
\end{choices}
\question{29. What does the underlined word "demoralizing" in paragraph 2 mean?}
\begin{choices}(2)
  \item Discouraging.
  \item Demanding.
  \item Comforting.
  \item Eye-opening.
\end{choices}
\question{30. Why was the speech composed in Harrell's pre-ALS voice?}
\begin{choices}(2)
  \item To enable his communication in real time.
  \item To make his speech emotionally meaningful.
  \item To improve the accuracy of speech translation.
  \item To help interpret his brain signals more easily.
\end{choices}
\question{31. What is Dr. John Ngai's attitude towards the BrainGate BCI?}
\begin{choices}(2)
  \item Favourable.
  \item Unclear.
  \item Cautious.
  \item Dismissive.
\end{choices}
\passagelabel{D}
\par It might surprise you to learn that dogs are playing an increasingly important role as conservationists. For centuries, they have been known as man's best friends. Now, their unique abilities are being made use of to help protect some of the world's most threatened species.\par
\par Quick and loyal, dogs possess many qualities that make them ideal for fieldwork, but it's their extraordinary sense of smell that makes them invaluable for conservation work. With around 220 million olfactory receptors(嗅觉受体) compared to just 5 million in humans, dogs can detect scents(气味) and analyze them with astonishing accuracy, devoting 40 times more of their brain to scent analysis than we do.\par
\par Thanks to this incredible ability, some dogs are now being trained to save at-risk animals by sniffing(嗅) out their droppings. At the University of Washington's Conservation Canines facility,20specially trained dogs have successfully tracked killer whale droppings, which only float briefly on the ocean's surface. Analyzing these samples enables scientists to monitor the health of killer whale populations, which are under significant stress from food shortages, pollution, and boat traffic.\par
\par Beyond tracking animal droppings, conservation dogs are also trained for live animal detection, further showcasing their abilities. They can locate rare species in their natural habitats. Additionally, they are now used to tackle the non-native species that threaten the survival of the local wildlife. Dogs also take on other important roles. For example, at Kenya's Mombasa port, dogs sniff out illegal wildlife products such as rhino horns and ivory. Across Africa, they work on the frontlines against illegal hunting, serving as a powerful security force enhancer.\par
\par While dogs may not be the ultimate solution to every conservation challenge, their incomparable abilities make them invaluable companions in the fight to protect our planet's biodiversity.\par
\question{32. Why are the figures mentioned in the second paragraph?}
\begin{choices}(2)
  \item To highlight the superiority of dogs' sense of smell.
  \item To suggest the amazing accuracy of dogs' analysis.
  \item To recognize the importance of conservation work.
  \item To show dogs' impressive achievements in fieldwork.
\end{choices}
\question{33. How do dogs help protect killer whales?}
\begin{choices}(2)
  \item By tracking their movements.
  \item By detecting their droppings
  \item By monitoring their populations.
  \item By analyzing the collected samples.
\end{choices}
\question{34. What can we learn from the example given in paragraph 4?}
\begin{choices}(2)
  \item Wildlife products are popular in African countries.
  \item Non-native species affect the survival of local wildlife.
  \item Dogs are used to detect and prevent illegal wildlife trade.
  \item Dogs can accurately locate the natural habitats of rare species.
\end{choices}
\question{35. What is the author's purpose in writing the text?}
\begin{choices}(2)
  \item To appeal to people to protect biodiversity.
  \item To report various environmental challenges.
  \item To introduce dogs' roles in conservation efforts.
  \item To prove the conventional claim of dogs' qualities.
\end{choices}
\sectitle{第二节（共5小题；每小题2.5分，满分12.5分）}
\par 根据短文内容，从短文后的选项中选出能填入空白处的最佳选项。选项中有两项为多余选项。\par
\par When others outperform us, should we feel threatened or inspired? Should we fear or respect those who are more capable, more skilled or more successful than us? \blanknum{36}\par
\par A study conducted by educational psychologists at Brown University explored how kindergartners in the U. S. and China tell stories about peers who perform better than others. \blanknum{37} Chinese children often expressed respect for high achievers and a desire to imitate their success. In contrast, U.S. children's stories typically focused on negative reactions, including envy and being left out.\par
\par The findings reflect the cultural narratives (叙述) surrounding achievement in each society. In American culture, the "self" is defined by qualities, and abilities, which are largely fixed. Thus, while the achievement of others reflects their high ability, it becomes a mirror to reveal one's own lack of ability. \blanknum{38} As a result, Chinese children tend to view others' achievements not as signs of natural superiority, but as opportunities for learning and self-improvement.\par
\par \blanknum{39} In China, students' overall performance may be publicly recognized to encourage success and motivate further improvement. Teachers even invite outstanding graduates to share their experiences and lessons with current students to promote learning. In the U. S., such phenomena are less common as they can be seen as threats to self-esteem.\par
\par Ultimately, our perceptions and reactions to others' success are profoundly shaped by cultural narratives surrounding achievement. \blanknum{40} These narratives, in other words, determine whether we feel threatened or inspired, fearful or respectful, annoyed or eager to copy others.\par
\begin{choices}(1)
  \item It can be a source of motivation.
  \item These cultural differences extend to educational practices.
  \item Education facilitates personal growth and social progress.
  \item They influence our motivations, relationships and actions.
  \item Both groups told detailed stories but showed clear differences in content.
  \item Different responses can lead to constructive or destructive reactions and actions.
  \item However, Chinese culture stresses abilities and qualities can improve through efforts.
\end{choices}
\parttitle{第三部分 语言知识运用（共两节，满分30分）}
\sectitle{第一节（共15小题；每小题1分，满分15分）}
\par 阅读下面短文，从短文后各题所给的 A、B、C 和 D 四个选项中，选出可以填入空白处的最佳选项。\par
\par Four-year-old Alan, who suffered from a serious disease, had one dream — to be a garbage truck driver. He was on the weekly lookout for trucks, and he even \blanknum{41} himself as a garbage man at the kindergarten costume show.\par
\par On July 10, Alan went one step further when a charity and his local waste management crew made him an honorary \blanknum{42} for three days. He had the opportunity to participate in the \blanknum{43} of the city beauticians.\par
\par At dawn, Alan started the road trip with the crew members. "Alan was a little \blanknum{44} when he first started, for early mornings were tough for him," said one of the Waste Management drivers. \blanknum{45}, that small reality check was not enough to reduce Alan's \blanknum{46}. After getting used to his early hours, he got to clock in with a special timecard and \blanknum{47} learn what it takes to be a garbage man. Alan got to \blanknum{48} every part of being a garbage truck man.\par
\par His day ended with a special \blanknum{49} as he got to ride alongside another driver. A crowd of community members gathered to \blanknum{50} Alan for what he had done. It was a truly \blanknum{51} day for him.\par
\par The amazing experience was all thanks to the local charity organization, and Alan wasn't the only person who got to \blanknum{52} the day. "We are so happy the supports from all the kind-hearted individuals always \blanknum{53} with him," Alan's mom commented. "We are so happy he finally \blanknum{54} his dream. The kindness of others has truly made a \blanknum{55} in Alan's life," she added with a grateful smile.\par
\cz{41}{classified}{dressed}{recognized}{treated}
\cz{42}{captain}{pilot}{citizen}{collector}
\cz{43}{routine}{strategy}{ceremony}{program}
\cz{44}{confident}{proud}{resistant}{tired}
\cz{45}{Besides}{However}{Therefore}{Otherwise}
\cz{46}{pressure}{enthusiasm}{burden}{responsibility}
\cz{47}{regularly}{immediately}{gradually}{patiently}
\cz{48}{experience}{control}{adjust}{promote}
\cz{49}{send-off}{check-in}{warm-up}{show-off}
\cz{50}{blame}{forgive}{thank}{punish}
\cz{51}{relaxing}{incredible}{challenging}{regretful}
\cz{52}{miss}{transform}{survive}{enjoy}
\cz{53}{stick}{agree}{compete}{negotiate}
\cz{54}{set aside}{worked out}{lived out}{went after}
\cz{55}{deal}{point}{choice}{difference}
\sectitle{第二节（共10小题；每小题1.5分，满分15分）}
\par 阅读下面短文，在空白处填入 1 个适当的单词或括号内单词的正确形式。\par
\par The origins of the umbrella can be traced back over thousands of years to early human societies. Initially used as sunshades, these early umbrellas \blanknum{56} (fashion) from materials like silk or paper, and waterproofed with wax(蜡) and lacquer (漆) \blanknum{57} (offer) rain protection.\par
\par According to Chinese legend, the first umbrella in ancient China was invented by Lady Yun, the wife of Lu Ban, a renowned craftsman in Chinese history. At the time, people relied on straw hats and capes during rainy days, \blanknum{58} proved inconvenient for travel. Inspired by \blanknum{59} sight of people taking shelter under eaves(屋檐), Lady Yun created a folding device made of bamboo and fabric. LuBan later improved this design, making it \blanknum{60} (practical) for rain and sun protection than before, paving the way \blanknum{61} the modern umbrella—a proof of the wisdom and \blanknum{62} (creative) of early Chinese people.\par
\par Pioneer \blanknum{63} China may be in the development of umbrellas, similar devices appeared in other ancient civilizations, such as Egypt, Assyria, and Greece, primarily for sun protection. The word "umbrella" \blanknum{64} (it) originates from the Latin "umbra", meaning "shade". Over time, as waterproofing techniques advanced, umbrellas changed from simple sunshades to essential tools for rain protection, \blanknum{65} (gain) widespread popularity across the globe.\par
\parttitle{第四部分 写作（共两节，满分 40 分）}
\sectitle{第一节（满分15分）}
\par 假定你是李华，外教 Chris 计划周一至周五在你班开展课外英语拓展阅读 (Supplementary Reading)。你对此活动很感兴趣，请你给他写一封邮件，内容包括：(1) 你的看法；(2) 你的建议，注意：\par
\par （1）写作词数应为 80 左右；\par
\par （2）请按如下格式在答题纸的相应位置作答。\par
\sectitle{第二节（满分25分）}
\par 阅读下面材料，根据其内容和所给段落开头语续写两段，使之构成一篇完整的短文，\par
\par A Game of Growth\par
\par The gym was alive with the sounds of squeaking (嘎吱作响) shoes, bouncing basketballs, and teammates shouting plays. It was game day, and I should have been focused, but my mind was elsewhere. My thoughts dragged me down, and I couldn't shake the heaviness I felt.\par
\par Every time I passed the ball, Ryan, one of our best players, would make some snide (挖苦的)remarks: "Can’t you aim better, Eliot?" or "Guess we know who won’t be taking the game-winning shot today." His words cut deep, feeding my doubts. I wanted to fire back, but the energy wasn’t there.\par
\par Coach Rivera noticed something was off. "Eliot, can I talk to you for a second?" he said after practice. His voice was calm but firm, the way it always was when he saw something wrong. I sighed, knowing I couldn't avoid the conversation. "Yeah, Coach?" He studied me for a moment, his eyes serious. "You're not yourself lately. I've also noticed there's tension between you and Ryan. What's going on?"\par
\par I hesitated, then said, "Nothing, I'm fine."\par
\par Coach didn't buy it. "I've been around long enough to know when something's weighing on a player. You're not just dealing with basketball issues, are you?"\par
\par I finally admitted, my throat tight and my eyes starting to water, "I have some family problems...my parents had an accident, and I've had to take care of them."\par
\par Coach nodded, understanding. "That's a lot to handle. But Ryan .. he's a tough guy, and you two are teammates, and the team only works if you're united. Let's talk things through."\par
\par The next day, Coach pulled Ryan and me aside. The conversation was awkward at first. Ryan explained that he felt I wasn't pulling my weight in the games. I told him about my situation, though I didn't go into details.\par
\par Ryan's expression softened. "I didn't know you were dealing with so much. I'm sorry for pushing you too hard," he said. "I could've handled it better," I admitted.\par
\par （1）续写词数应为 150 左右；\par
\par （2）请按如下格式在答题卡的相应位置作答。\par
\writingpara{Paragraph 1: With Coach's help, Ryan and I started working together instead of against each other: \blanknum{}}
\writinglines{8}
\writingpara{Paragraph 2: In the championship game, Ryan passed me the ball for the game-winning shot.}
\writinglines{8}
\end{document}
